\documentclass[12pt]{article}
\usepackage[utf8]{inputenc}
\usepackage[T1]{fontenc}
\usepackage{hyperref}
\title{Guinea-Bissau Creole Language Resources for K-12 Curriculum Development}
\author{Compiled by Claude Sonnet 5.5 (claude-sonnet-5-5) for Andr\'as Kornai}
\date{}
\begin{document}
\maketitle

\section{Introduction}
Guinea-Bissau Creole (Kriol, Kriyol, crioulo guineense; ISO 639-3 \texttt{pov}) is a Portuguese-based creole and the lingua franca of Guinea-Bissau. Most of what this search turned up is academic work, mostly in Portuguese; it is not exhaustive (one search).

\section{Primary Educational Resources}
No official school textbooks or primers for teaching Kriol were found. Academic studies argue for teaching Kriol in schools:
\begin{itemize}
\item Vladimir Renato Nanfadi, thesis (UNILAB) on Kriol in school education. \url{https://repositorio.unilab.edu.br/jspui/bitstream/123456789/6735/1/Vladimir%20Renato%20Nanfadi.pdf}
\item Article in \emph{Revista Mandinga} (UNILAB) on language and schooling. \url{https://revistas.unilab.edu.br/index.php/mandinga/article/download/34/24/138}
\item Survey of 25 Guineans on language policy in school (UFSC, \emph{F\'orum}). \url{https://periodicos.ufsc.br/index.php/forum/article/download/95286/59081/411214}
\end{itemize}

\section{Linguistic and Reference Materials}
\begin{itemize}
\item Cirineu Cecote Stein, bilingual Portuguese--Guinea-Bissau Creole dictionary (vol.\ II), Editora UFPB, 2021. \url{https://www.editora.ufpb.br/sistema/press5/index.php/UFPB/catalog/download/705/941/8096?inline=1}
\item Scantamburlo, article on the official orthography of Guinea-Bissau Creole. \url{https://repositorio.ensinolusofona.pt/bitstream/10437.1/12021/1/S4_2018_Scantamburlo_Grafia%20oficial%20do%20crioulo%20guineense.pdf}
\item Alain Kihm, \emph{Kriyol Syntax} (John Benjamins, 1994). \url{https://www.degruyterbrill.com/document/doi/10.1075/cll.14/html?lang=en}
\item Survey chapter: Guinea-Bissau Kriyol, Atlas of Pidgin and Creole Language Structures. \url{https://apics-online.info/surveys/33}
\end{itemize}

\section{Religious and Community Resources}
None found in this search.

\section{Digital and Audio Resources}
None found in this search.

\section{Curriculum Development Context}
Portuguese is the official language and school language; Kriol has been disadvantaged in school, and sources report divergent views on making it official. Scantamburlo's \emph{Dicion\'ario Guineense-Portugu\^es} (FASPEBI, Bissau) is mentioned in the results without a link. A Kriol curriculum would need to begin with the dictionary, orthography and grammar above and create classroom materials.
\end{document}